\chapter{Υλοποίηση σε \tl{CPU}}
\label{ch:cpu}

\section{Σχεδιαστικές επιλογές}
\label{sec:cpu-design}

\subsection{Διανυσματοποίηση}
\label{sec:cpu-simd}

\subsection{Αφαίρεση διανυσματικών τύπων}
\label{sec:cpu-simd-traits}

\subsection{Υποστήριξη αριθμών κινητής υποδιαστολής μισής ακρίβειας}
\label{sec:cpu-fp16}

\subsection{Παραλληλοποίηση σε πολλαπλά νήματα}
\label{sec:cpu-threads}

\section{Διτονική ταξινόμηση}
\label{sec:cpu-bitonic}

Το χρονοδιάγραμμα του περικομμένου δικτύου κατασκευάζεται μία φορά και είναι
κοινό και για τις τρεις αρχιτεκτονικές. Αποτελείται από δύο φάσεις. Στην πρώτη
εκτελείται κανονική διτονική ταξινόμηση μέχρι να σχηματιστούν ταξινομημένα
μπλοκ μεγέθους $k'$, όπου $k'$ η μικρότερη δύναμη του δύο που δεν είναι
μικρότερη του ζητούμενου \tl{k}. Στη δεύτερη, τα μπλοκ συγχωνεύονται ανά ζεύγη
και κάθε συγχώνευση κρατά μόνο τα $k'$ καλύτερα στοιχεία των $2k'$ που
συγκρίθηκαν, οπότε το μέγεθος του πίνακα υποδιπλασιάζεται σε κάθε βήμα. Η
διαδικασία επαναλαμβάνεται έως ότου απομείνει ένα μόνο μπλοκ.

Η άμεση μεταφορά αυτού του χρονοδιαγράμματος σε κώδικα οδηγεί σε υλοποίηση
απελπιστικά περιορισμένη από τη μνήμη. Για \tl{k} ίσο με 256, η πρώτη φάση και
μόνο αποτελείται από 36 επίπεδα, καθένα από τα οποία θα ήταν ένα ξεχωριστό
πέρασμα πάνω σε ολόκληρο τον πίνακα. Σε είσοδο τεσσάρων εκατομμυρίων στοιχείων
αυτό σημαίνει δεκάδες διαδοχικές αναγνώσεις και εγγραφές δεκαέξι
\tl{MB}, με ελάχιστο υπολογισμό ανάμεσά τους. Ολόκληρη η υλοποίηση της
παρούσας εργασίας είναι οργανωμένη γύρω από την αποφυγή αυτού ακριβώς του
σχήματος.

Η βασική τεχνική είναι η συγχώνευση διαδοχικών επιπέδων σε ένα πέρασμα. Ένα
επίπεδο χαρακτηρίζεται από την απόσταση $j$ μεταξύ των στοιχείων που
συγκρίνονται, και η απόσταση αυτή καθορίζει τι είναι εφικτό. Όταν το $j$ είναι
μικρότερο από το πλάτος του διανυσματικού καταχωρητή, τα δύο στοιχεία κάθε
σύγκρισης βρίσκονται ήδη μέσα στον ίδιο καταχωρητή. Mια ολόκληρη ακολουθία
τέτοιων επιπέδων εκτελείται τότε ως μία φόρτωση, μια σειρά αναδιατάξεων και
συγκρίσεων εντός καταχωρητών, και μία αποθήκευση. Όταν το $j$ είναι μεγαλύτερο
αλλά το σύνολο εργασίας ενός τμήματος χωρά στην κρυφή μνήμη, η ακολουθία
εκτελείται τμήμα προς τμήμα, με το κάθε τμήμα να παραμένει στην κρυφή μνήμη για
όλα τα επίπεδα. Μόνο τα επίπεδα με πολύ μεγάλο $j$ απαιτούν ξεχωριστό πέρασμα.
Επιπλέον, ένα επίπεδο περικοπής του οποίου η απόσταση ισούται με το πλάτος του
καταχωρητή απορροφά και τα επίπεδα επαναταξινόμησης που το ακολουθούν.

Το αποτέλεσμα είναι ότι η κίνηση δεδομένων δεν είναι απλώς μετρήσιμη, όπως
περιγράφηκε στην ενότητα~\ref{sec:meth-traffic}, αλλά και προβλέψιμη από τη δομή
του χρονοδιαγράμματος. Η πρώτη φάση συμπτύσσεται σε ένα πέρασμα επί τόπου,
δηλαδή μία ανάγνωση και μία εγγραφή ανά στοιχείο, άρα $2n$. Κάθε βήμα της
δεύτερης φάσης διαβάζει $n_i$ στοιχεία και γράφει $n_i/2$, και η
επαναταξινόμηση που ακολουθεί προσθέτει άλλα $2 \cdot n_i/2$, συνολικά
$2{,}5\,n_i$. Επειδή το $n_i$ υποδιπλασιάζεται σε κάθε βήμα, το άθροισμα της
δεύτερης φάσης τείνει στο $5n$. Η συνολική κίνηση είναι επομένως περίπου $7n$
στοιχεία, που είναι η προέλευση του συντελεστή επανάληψης της
ενότητας~\ref{sec:meth-traffic}. Για $n = 2^{22}$ και \tl{k} ίσο με 256 η
έκφραση αυτή δίνει 117\,435\,392 \tl{bytes}, τιμή που ταυτίζεται ακριβώς με τη
μετρούμενη.

Η ταυτόχρονη εκτέλεση σε πολλά νήματα γίνεται εντός κάθε ομάδας και όχι μεταξύ
ομάδων, καθώς κάθε ομάδα εξαρτάται από το αποτέλεσμα της προηγούμενης. Το εύρος
εξόδου κάθε ομάδας μοιράζεται στα νήματα, με τα όρια ευθυγραμμισμένα στο μέγεθος
του τμήματος ή της γραμμής κρυφής μνήμης ώστε δύο νήματα να μη γράφουν ποτέ στην
ίδια γραμμή. Ανάμεσα στις ομάδες μεσολαβεί φράγμα συγχρονισμού. Το πλήθος των
νημάτων περιορίζεται επιπλέον ανάλογα με το μέγεθος της εισόδου, ώστε για μικρές
εισόδους το κόστος του συγχρονισμού να μην υπερβαίνει το όφελος της
παραλληλίας.

Τα επίπεδα περικοπής δεν μπορούν να εκτελεστούν επί τόπου, διότι κάθε στοιχείο
της εξόδου προκύπτει από δύο στοιχεία της εισόδου σε διαφορετικές θέσεις. Για
τον λόγο αυτό διατηρείται ένας δεύτερος πίνακας ίδιου μεγέθους και οι δύο
πίνακες εναλλάσσονται μετά από κάθε τέτοια ομάδα. Η επιπλέον μνήμη είναι το
τίμημα για την αποφυγή ενδιάμεσων αντιγραφών.

Τέλος, κάθε διανυσματική διαδρομή συνοδεύεται από ισοδύναμη βαθμωτή. Αν το υλικό
δεν υποστηρίζει το απαιτούμενο πλάτος, ή αν οι παράμετροι ενός επιπέδου
βρίσκονται εκτός του εύρους που χειρίζεται η διανυσματική υλοποίηση, η ίδια
ακολουθία επιπέδων εκτελείται βαθμωτά. Το αποτέλεσμα είναι ταυτόσημο και
διαφέρει μόνο ο χρόνος, κάτι που επιτρέπει στον έλεγχο ορθότητας της
ενότητας~\ref{sec:meth-correctness} να επικυρώνει και τις δύο διαδρομές.

\section{\tl{Map-reduce} με τοπικούς σωρούς}
\label{sec:cpu-mapreduce}

\section{Διαχείριση μνήμης και κίνηση δεδομένων}
\label{sec:cpu-memory}

\section{Υλοποίηση αναφοράς}
\label{sec:cpu-reference}

\section{Περιοριστικοί παράγοντες}
\label{sec:cpu-bottlenecks}
